\section{Evaluation}

The evaluation asks whether the implementation can support controlled behavior
studies from hardware trace records. We evaluate trace correctness, source
labeling, code attribution, and reproducibility. Board transport throughput,
cycle overhead, and RAM boot are outside the evaluated scope.

\begin{table*}[t]
\centering
\caption{Evaluation questions, supporting evidence, and boundaries.}
\label{tab:evaluation}
\scriptsize
\begin{tabularx}{\textwidth}{p{0.08\textwidth}p{0.24\textwidth}p{0.20\textwidth}p{0.17\textwidth}Y}
\toprule
RQ & Claim tested & Evidence & Result & Boundary \\
\midrule
RQ1 & Trace records preserve ordering and syscall pairing needed for reconstruction & Directed trace tests & Covered by the current test corpus & Does not measure production transport throughput or runtime slowdown \\
RQ2 & Reconstructed fields identify whether they came from hardware trace, ELF metadata, runtime maps, or reference logs & Reconstruction summaries & All reported fields carry a source label & Reference-derived fields validate expectations but are not hardware recovery \\
RQ3 & Code attribution handles common mapping edge cases & Local attribution tests & Attribution edge cases are covered locally & Local fixtures do not establish board-native source-line recovery \\
RQ4 & The evidence package is reproducible and auditable & Local and Docker runs & Quick, local, full, and Docker reproduction complete & Reproduction covers the recorded artifact, not experiments outside the manifest \\
RQ5 & Controlled workloads exercise behavior reconstruction without becoming a malware benchmark & Case-study manifests and behavior checks & Controlled behavior checks are recorded in the artifact package & Does not imply real-malware accuracy, TTP coverage, or payload equivalence \\
\bottomrule
\end{tabularx}
\end{table*}


\subsection{RQ1: Trace Correctness}

The directed tests check event ordering and syscall pairing at the trace-record
level. They cover trap/retire separation, strict syscall entry/return pairing,
same-cycle ordering, dual-commit ordering, privilege transitions, and negative
relaxed-SRET cases.

\subsection{RQ2: Source Labeling}

The source-labeling checks verify that reconstructed fields identify their
source: hardware trace, ELF metadata, runtime map, or reference log. Reference
traces can validate expected behavior, but they are not reported as
hardware-recovered fields.

\subsection{RQ3: Code Attribution}

The local code-analysis tests cover ELF identity, PIE/ASLR load bias, dynamic
object attribution, fork/exec ownership, stripped binaries, and source-map
boundaries. These tests support attribution logic, but they are not reported as
a new board run or as board-native DWARF evidence.

\subsection{RQ4: Reproducibility}

The artifact package is organized around a canonical evidence directory and
machine-readable manifests. Quick local reproduction is:

\begin{verbatim}
uv run rvmt repro:quick
\end{verbatim}

Docker-local reproduction is:

\begin{verbatim}
uv run rvmt ndss:docker-full
\end{verbatim}

The current project status records successful quick, local, full, and
Docker-local reproduction paths. Experiments that were not run are recorded
separately and are not included in the results.
